% Build from this directory: latexmk -pdf report.tex
\documentclass[12pt,letterpaper]{article}

% Page layout and typography (pdfLaTeX).
\usepackage[T1]{fontenc}
\usepackage[utf8]{inputenc}
\usepackage{lmodern} % Scalable fonts required by microtype font expansion.
\usepackage[margin=1in]{geometry}
\usepackage{microtype}
\usepackage{amsmath,amssymb}
\usepackage{graphicx}
\graphicspath{{figures/}{../results/figures/}}
\usepackage{booktabs,tabularx}
\usepackage{caption}
\usepackage[numbers,sort&compress]{natbib}
\usepackage[hidelinks]{hyperref}
\urlstyle{same}
\setlength{\emergencystretch}{2em}

% Edit these fields before submitting.
\newcommand{\reporttitle}{FP-Growth for Frequent Itemset Mining}
\newcommand{\reportauthor}{Your Name}
\newcommand{\reportcourse}{Course Number and Name}
\newcommand{\reportinstructor}{Instructor Name}
\newcommand{\reportdate}{October 30, 2026}
\hypersetup{pdftitle={\reporttitle},pdfauthor={\reportauthor}}

% Visible writing prompts. Replace every occurrence before submission.
\newcommand{\placeholder}[1]{\par\noindent\textit{[To write: #1]}\par}

\begin{document}
\hypersetup{pageanchor=false}

% Title Page: 

% \begin{titlepage}
%   \centering
%   {\large Stony Brook University\par}
%   \vspace{2cm}
%   {\LARGE\bfseries \reporttitle\par}
%   \vspace{0.5cm}
%   {\large Implementation, Validation, and Performance Evaluation\par}
%   \vspace{2cm}
%   {\large \reportauthor\par}
%   \vspace{1cm}
%   {\large \reportcourse\par}
%   \vspace{0.5cm}
%   {\large Instructor: \reportinstructor\par}
%   \vfill
%   {\large \reportdate\par}
% \end{titlepage}

\pagenumbering{roman}
\hypersetup{pageanchor=true}
\begin{abstract}
\placeholder{Summarize the problem, implementation, datasets, validation
approach, and main measured findings in about 150--250 words. Write this last.}
\end{abstract}
\tableofcontents
% Add \listoffigures or \listoftables only if the final report needs them.
\clearpage
\pagenumbering{arabic}

\section{Introduction}
\label{sec:introduction}
\placeholder{Introduce frequent itemset mining, explain the project goals,
and summarize the implementation and evaluation contributions.}
% Example citation syntax: FP-growth was introduced by \citet{han2000fpgrowth}.

\section{Background and Algorithm}
\label{sec:algorithm}
\subsection{Problem Definition}
\placeholder{Define transactions, items, itemsets, absolute and relative
support, and the minimum support threshold. State how a fractional threshold
is converted to a transaction count.}
% Example equation:
% \begin{equation}
%   \operatorname{supp}(X) = |\{T \in D : X \subseteq T\}|.
%   \label{eq:support}
% \end{equation}

\subsection{FP-Tree Construction}
\placeholder{Explain frequency counting, removal of infrequent items,
deterministic item ordering, prefix sharing, and the header table.
Cite the original algorithm using \texttt{\string\citep\{han2000fpgrowth\}}.}

\subsection{Recursive Pattern Mining}
\placeholder{Explain conditional pattern bases, conditional FP-trees,
recursion, stopping conditions, and any single-path optimization implemented.
Include pseudocode or a small tree diagram where helpful.}

\subsection{Computational Considerations}
\placeholder{Discuss tree size, recursion, output size, and factors that
affect runtime and memory. Distinguish theoretical expectations from measurements.}

\section{Implementation}
\label{sec:implementation}
\subsection{Program Structure and Data Structures}
\placeholder{Describe the language, main modules, node representation,
header links, support counts, memory ownership, and output format.}

\subsection{Input Handling and Design Choices}
\placeholder{Document input assumptions, duplicate-item handling,
tie breaking, parameter validation, and any implementation optimizations.}

% Optional bonus section: uncomment only if a CUDA version is implemented.
% \subsection{CUDA Implementation}
% Describe CPU/GPU work division, kernels, memory transfers, synchronization,
% device limits, and how GPU results are checked against the CPU version.

\section{Datasets and Preprocessing}
\label{sec:datasets}
\subsection{Benchmark Data}
\placeholder{Describe the UCI Gas Sensor Array Drift Dataset and any other
benchmarks actually used. Cite the dataset with
\texttt{\string\citep\{vergara2012gas\}}. Record transaction counts, item
counts, transaction lengths, and the selected batches or subsets.}

\subsection{Transaction Encoding}
\placeholder{Document selected features, normalization, binning rules,
threshold fitting, item-ID mapping, and treatment of gas labels and batches.
Explain how continuous measurements become discrete transactions.}

\section{Correctness Validation}
\label{sec:validation}
\subsection{Worked Example}
\placeholder{Show a small transaction database and minimum support.
List the expected frequent itemsets and exact support counts, illustrate
the FP-tree and at least one conditional tree, and compare actual output.}

\subsection{Additional Tests}
\placeholder{Report the tests actually run and their outcomes. Consider
empty input, repeated items, identical transactions, support boundaries,
no frequent items, tied frequencies, and comparison with exhaustive enumeration
on small databases.}

\section{Experimental Methodology}
\label{sec:methodology}
\subsection{Hardware and Software}
\placeholder{Record CPU, RAM, operating system, compiler version and flags,
and GPU/CUDA details if applicable. Include commands needed to reproduce runs.}

\subsection{Measurements and Experimental Parameters}
\placeholder{Specify support thresholds, input sizes, repetitions, warm-up
policy, timer, and memory measurement method. Define whether timings include
preprocessing, input/output, tree construction, mining, and GPU transfers.
State how variability and incomplete runs are reported.}

\section{Results and Discussion}
\label{sec:results}
\subsection{Execution Time}
\placeholder{Present measured runtime against minimum support and input
size, with units and variability. Explain observed trends.}

\subsection{Memory Usage and Output Size}
\placeholder{Report peak memory and number of frequent itemsets for the
same configurations. Explain the relationship between support, tree size,
output size, and resource use.}

\subsection{Limitations}
\placeholder{Discuss dataset and preprocessing choices, measurement limits,
scalability, and any failed or incomplete experiments.}

% Copy this table when measurements are available; do not invent values.
% \begin{table}[htbp]
%   \centering
%   \caption{Performance on the specified benchmark configuration.}
%   \label{tab:performance}
%   \begin{tabular}{rrrr}
%     \toprule
%     Min. support (\%) & Itemsets & Time (s) & Peak memory (MiB) \\
%     \midrule
%     % Insert measured rows here.
%     \bottomrule
%   \end{tabular}
% \end{table}
%
% \begin{figure}[htbp]
%   \centering
%   \includegraphics[width=0.85\linewidth]{runtime.pdf}
%   \caption{Runtime versus minimum support; specify dataset and run settings.}
%   \label{fig:runtime}
% \end{figure}
% Cross-reference with Figure~\ref{fig:runtime} or Table~\ref{tab:performance}.

\section{Conclusion}
\label{sec:conclusion}
\placeholder{Summarize what was implemented, the evidence for correctness,
the main measured results, and specific future improvements.}

\clearpage
\phantomsection
\addcontentsline{toc}{section}{References}
% Keep starter references visible while drafting. Remove this \nocite line
% once real citations have been added; then only cited entries will appear.
\nocite{han2000fpgrowth,vergara2012gas}
\bibliographystyle{plainnat}
\bibliography{references}

\appendix
\section{Reproducibility Details}
\label{app:reproduction}
\placeholder{Include exact build and run commands, preprocessing parameters,
and additional validation or result tables that would interrupt the main text.
Remove this appendix if all details fit in the main report.}

\end{document}
